\documentclass[11pt]{article}

% ---- Portable preamble (compiles with XeLaTeX or pdfLaTeX, no project macros) ----
\usepackage{amsmath,amsthm,amssymb}
\usepackage[margin=1in]{geometry}
\usepackage[colorlinks=true,linkcolor=blue,urlcolor=blue,citecolor=blue]{hyperref}
\usepackage{mathtools}
\usepackage{enumitem}
\usepackage{booktabs}
\usepackage{array}
\usepackage[round,authoryear]{natbib}

\theoremstyle{plain}
\newtheorem{theorem}{Theorem}
\newtheorem{lemma}{Lemma}
\newtheorem{corollary}{Corollary}
\newtheorem{proposition}{Proposition}
\theoremstyle{definition}
\newtheorem{assumption}{Assumption}
\newtheorem{definition}{Definition}
\theoremstyle{remark}
\newtheorem{remark}{Remark}

\newcommand{\E}{\mathbb{E}}
\newcommand{\PP}{\mathbb{P}}
\newcommand{\Pn}{\mathbb{P}_n}
\newcommand{\R}{\mathbb{R}}
\newcommand{\cX}{\mathcal{X}}
\newcommand{\cT}{\mathcal{T}}
\newcommand{\cS}{\mathcal{S}}
\newcommand{\cL}{\mathcal{L}}
\newcommand{\ez}{e_0}
\newcommand{\mz}{\mu_0}
\newcommand{\wz}{w_0}
\newcommand{\ehat}{\widehat{e}}
\newcommand{\mhat}{\widehat{\mu}}
\newcommand{\what}{\widehat{w}}
\newcommand{\that}{\widehat{\tau}}
\newcommand{\thetahat}{\widehat{\theta}}
\newcommand{\pihat}{\widehat{\pi}}
\newcommand{\Vhat}{\widehat{V}}
\newcommand{\dhat}{\widehat{\delta}}
\newcommand{\Ltwo}[1]{\lVert #1 \rVert_{2}}
\newcommand{\Lw}[1]{\lVert #1 \rVert_{2,w}}
\DeclareMathOperator*{\argmin}{arg\,min}
\DeclareMathOperator{\Var}{Var}
\DeclareMathOperator{\Cov}{Cov}
\DeclarePairedDelimiter{\abs}{\lvert}{\rvert}
\DeclarePairedDelimiter{\card}{\lvert}{\rvert}

\title{Interpretable Causal Inference with Optimal Decision Trees\\
\large Notation, assumptions, and main results}
\author{}
\date{2026-08-06}

\begin{document}
\maketitle

\begin{abstract}
This note fixes notation for the whole paper and states the three main results
without proofs, so that the symbol system and the assumption set can be checked
before drafting. Three substantive changes from earlier drafts are flagged
in-line and collected in Section~\ref{sec:changes}: positivity is weakened to a
\emph{one-sided} condition, the rate requirement is stated in \emph{product}
form, and the loss--norm constant is corrected. Section~\ref{sec:notation}
resolves five symbol collisions that were live across the previous drafts.
\end{abstract}

\tableofcontents

%======================================================================
\section{Notation}\label{sec:notation}
%======================================================================

Every symbol used anywhere in the paper is defined here. Symbols are grouped by
role, and Section~\ref{sec:collisions} records the collisions this system
resolves --- five symbols previously carried two or three distinct meanings
across the draft set.

\subsection{Data, target parameter, and the score}

\begin{center}
\begin{tabular}{@{}ll@{\hspace{2em}}l@{}}
\toprule
Symbol & Meaning & Notes \\
\midrule
$O=(X,A,Y)$ & one observation & $O_1,\dots,O_n$ i.i.d.\ from $P$ \\
$n$ & sample size & \\
$X\in\cX$ & covariate vector & \\
$d$ & dimension of $X$ & \textbf{not} $p$; see Sec.~\ref{sec:collisions} \\
$A\in\{0,1\}$ & binary treatment & \\
$Y\in\R$ & observed outcome & \\
$Y(a)$ & potential outcome under $A=a$ & $a\in\{0,1\}$ \\
$P$ & distribution of $O$ & \\
$\Pn$ & empirical measure & $\Pn f=n^{-1}\sum_i f(O_i)$ \\
$\pi$ & $\PP(A=1)$ & $\pihat=\Pn A$ \\
$\theta_0$ & the ATT, $\E[Y(1)-Y(0)\mid A=1]$ & target parameter \\
$\psi$ & efficient influence function & Eq.~\eqref{eq:score} \\
$V$ & asymptotic variance of $\thetahat$ & Hahn ATT efficiency bound \\
$\Vhat$ & its estimator & Eq.~\eqref{eq:vhat} \\
$z_{1-\alpha/2}$ & standard normal quantile & \\
\bottomrule
\end{tabular}
\end{center}

\subsection{Covariate space}

The paper works in two tiers. The \emph{discrete tier} is primary; the
\emph{continuous tier} is an explicitly scoped bridge.

\begin{center}
\begin{tabular}{@{}ll@{\hspace{2em}}l@{}}
\toprule
Symbol & Meaning & Tier \\
\midrule
$\cX=\prod_{j=1}^d\cX_j$ & finite product covariate space & discrete \\
$J$ & number of levels per coordinate & discrete; \textbf{not} $K$ \\
$M=\card{\cX}$ & total number of cells, $M\le J^d$ & discrete \\
$p_x=\PP(X=x)$ & probability of cell $x$ & discrete \\
$s$ & \emph{local sparsity}: number of active coordinates & both \\
$\bar\alpha$ & harmonic-mean H\"older smoothness & continuous only \\
$K$ & number of cross-fitting folds & Part~I only \\
$n_k$ & size of fold $k$ & Part~I only \\
\bottomrule
\end{tabular}
\end{center}

\begin{remark}[$s$ is local, and $d$ is not constrained by it]
$s$ counts active coordinates \emph{within a region}. The active set may differ
across regions, so $\gamma_0$ may depend on all $d$ coordinates globally while
$s$ remains small. This is the property that makes trees the right model class
here and is developed at Remark~\ref{rem:additive}.
\end{remark}

\subsection{Nuisance functions}

\begin{center}
\begin{tabular}{@{}ll@{\hspace{2em}}l@{}}
\toprule
Symbol & Meaning & Notes \\
\midrule
$\ez(x)$ & $\PP(A=1\mid X=x)$, propensity & \\
$\mz(x)$ & $\E[Y\mid A=0,X=x]$, control outcome & controls only \\
$\gamma_0$ & generic true nuisance & $\gamma_{0,e}=\ez$, $\gamma_{0,\mu}=\mz$ \\
$j$ & nuisance index & $j\in\{e,\mu\}$ \\
$\eta=(\ehat,\mhat)$ & the nuisance pair & \\
$\wz(x)$ & ATT weight $\ez(x)/\{1-\ez(x)\}$ & $\what$ its estimate \\
$c$ & positivity constant & $1-\ez(x)\ge c$; see A\ref{ass:causal} \\
$\rho(\cdot,\cdot)$ & loss function & \textbf{not} $L$; see Sec.~\ref{sec:collisions} \\
$C_{\mathrm{ln}}$ & loss--norm link constant & $=1$ for squared error \\
\bottomrule
\end{tabular}
\end{center}

\subsection{Trees, partitions, and selection}

\begin{center}
\begin{tabular}{@{}ll@{\hspace{2em}}l@{}}
\toprule
Symbol & Meaning & Notes \\
\midrule
$\tau$ & a tree partition of $\cX$ & recursive binary, axis-aligned \\
$\ell$ & a leaf of $\tau$ & \\
$\card{\tau}$ & number of leaves of $\tau$ & \\
$\bar L$ & leaf budget & $L_n$ when it grows with $n$ \\
$\cT_{\bar L}$ & tree partitions with $\card{\tau}\le\bar L$ & finite set in the discrete tier \\
$\tau^*_j$ & minimal partition on which $\gamma_{0,j}$ is constant & \\
$L^*_j=\card{\tau^*_j}$ & its cardinality & \\
$\cS_j$ & sufficient class, $\{\tau\in\cT_{\bar L}:\tau$ refines $\tau^*_j\}$ & \\
$\that_j$ & selected partition for nuisance $j$ & data-adaptive \\
$m_n$ & minimum leaf count & \textbf{not} fold size \\
$\lambda_n$ & leaf penalty & $\lambda_n\to0$ \\
$R^{(j)}(\tau)$ & population risk of best predictor on $\tau$ & $R^{(j)}_n$ empirical \\
$\Delta_j$ & structural margin & Lemma~\ref{lem:margin} \\
\bottomrule
\end{tabular}
\end{center}

Structure selection is
\begin{equation}\label{eq:select}
  \that_j \;\in\; \argmin_{\tau\in\cT_{\bar L}}
    \Bigl\{ R^{(j)}_n(\tau) + \lambda_n\card{\tau} \Bigr\},
\end{equation}
a \emph{global} minimum over $\cT_{\bar L}$, and leaf values are the
within-leaf empirical means: for the propensity, the treated fraction in each
leaf; for the outcome, the control-unit mean in each leaf.

\subsection{Approximation error}

\begin{center}
\begin{tabular}{@{}ll@{\hspace{2em}}l@{}}
\toprule
Symbol & Meaning & Notes \\
\midrule
$\delta_j$ & $L_2$ approximation error of the best $\bar L$-leaf tree & Def.~\ref{def:delta} \\
$\delta_j(L)$ & the same as a function of the budget & \\
$q_j$ & decay exponent, $\delta_j(L)\le C_jL^{-q_j}$ & Part~III \\
$\dhat_j$ & estimator of $\delta_j$ & Prop.~\ref{prop:estimable} \\
$\delta_{\mathrm{fid}}$ & \emph{fidelity diagnostic} $\thetahat-\thetahat_{\mathrm{cf}}$ & \textbf{not} $\delta$ \\
$\Lw{\cdot}$ & $(1-\ez)$-weighted $L_2$ norm & Def.~\ref{def:delta} \\
\bottomrule
\end{tabular}
\end{center}

\begin{remark}[$\delta$ and $\delta_{\mathrm{fid}}$ are different objects]
\label{rem:two-deltas}
This distinction is easy to lose and was lost in the software. $\delta_j$ is an
\emph{approximation error} of a nuisance function and bounds a \emph{bias}.
$\delta_{\mathrm{fid}}=\thetahat-\thetahat_{\mathrm{cf}}$ is a \emph{difference
between two estimators} and bounds \emph{selection variance} relative to a
cross-fitted twin. They enter different intervals with different justifications.
The implementation currently calls the second one \texttt{delta}; it must be
renamed before any interval built on $\delta_j$ is implemented.
\end{remark}

\subsection{Symbol collisions this system resolves}\label{sec:collisions}

Each row was live in at least two files of the previous draft set.

\begin{center}
\begin{tabular}{@{}lp{6.2cm}l@{}}
\toprule
Symbol & Competing meanings & Resolution \\
\midrule
$s$ & sparsity (manuscript); H\"older smoothness (theory
      correspondence); leaf count (VC-dimension fragments) &
      $s$ = sparsity; $\bar\alpha$ = smoothness; $L$ = leaves \\
$K$ & cross-fitting folds (Part~I); levels per coordinate
      (entropy term $L\log(dK)$, Part~II) &
      $K$ = folds; $J$ = levels \\
$p$ & covariate dimension ($\cX\subset\R^p$); cell probability
      $p_x$; weak-$\ell_p$ index &
      $d$ = dimension; $p_x$ = cell probability; weak-$\ell_r$ \\
$m$ & minimum leaf count; fold sample size (in $\lambda_m\asymp\log(md)/m$) &
      $m_n$ = min leaf count; $n_k$ = fold size; $\lambda_n$ = penalty \\
$L$ & leaf count / budget; loss function $L\{Z,\gamma(X)\}$ &
      $L$ = leaves; $\rho(\cdot,\cdot)$ = loss \\
$\delta$ & nuisance approximation error; estimator-difference
      fidelity diagnostic &
      $\delta_j$ = approximation error; $\delta_{\mathrm{fid}}$ = diagnostic \\
\bottomrule
\end{tabular}
\end{center}

%======================================================================
\section{The estimator}
%======================================================================

The efficient influence function for the ATT \citep{hahn1998} is
\begin{equation}\label{eq:score}
  \psi(O;\theta,\eta)
  \;=\; \frac1\pi\Bigl[\,
    A\bigl\{Y-\mu(X)-\theta\bigr\}
    \;-\; (1-A)\,w(X)\bigl\{Y-\mu(X)\bigr\}
  \Bigr],
  \qquad w=\frac{e}{1-e}.
\end{equation}
Both terms have mean zero at $\eta=\eta_0$, $\theta=\theta_0$. Solving
$\Pn\psi(\cdot;\theta,\widehat\eta)=0$ gives the estimator
\begin{equation}\label{eq:att-estimator}
  \thetahat \;=\; \frac{1}{n\pihat}\sum_{i=1}^n
   \Bigl[\, A_i\bigl\{Y_i-\mhat(X_i)\bigr\}
   \;-\; (1-A_i)\,\what(X_i)\bigl\{Y_i-\mhat(X_i)\bigr\}\Bigr].
\end{equation}

\begin{remark}[Sign correction]\label{rem:sign}
The control-arm term carries a \textbf{minus}. Earlier drafts displayed a plus
while defining the score as in \eqref{eq:score}, so the displayed estimator was
not the one the score defines. Verified by direct solution of
$\Pn\psi=0$.
\end{remark}

The exact second-order remainder is bilinear, which is the structural fact the
whole paper rests on:
\begin{equation}\label{eq:bilinear}
  \E\bigl[\psi(O;\theta_0,\eta)\bigr]
  \;=\; \frac1\pi\,
  \E\Bigl[ \bigl\{1-\ez(X)\bigr\}\bigl\{\wz(X)-w(X)\bigr\}
           \bigl\{\mz(X)-\mu(X)\bigr\}\Bigr].
\end{equation}

\begin{remark}[This identity is standard; cite it, do not claim it]
Equation~\eqref{eq:bilinear} is the textbook doubly-robust computation
\citep{robins1995,vanderlaan2003,kennedy2024}. It is currently a standalone
seven-page note imported as ``Lemma~1'', i.e.\ presented as ours. It should
appear as a cited fact with the one-line derivation in an appendix.
\end{remark}

Two consequences of \eqref{eq:bilinear} that drive later definitions. First, the
weight $\{1-\ez\}$ sits on the $\mu$ factor, so the relevant norm for $\mu$ is
$(1-\ez)$-weighted --- this is why Definition~\ref{def:delta} is stated that way.
Second, $\abs{\wz-w}\le\abs{\ez-e}/\{(1-e)(1-\ez)\}$, so bounding
\eqref{eq:bilinear} needs $1-\ez$ and $1-\ehat$ bounded \emph{below} and nothing
else --- which is the basis for the positivity change below.

%======================================================================
\section{Assumptions}
%======================================================================

\subsection{Common to all parts}

\begin{assumption}[Identification and one-sided positivity]\label{ass:causal}
Consistency; $\{Y(0),Y(1)\}\perp A\mid X$; $\pi>0$; and there is $c\in(0,1]$
with
\[
  1-\ez(x)\;\ge\;c \qquad\text{for every }x\in\cX .
\]
Fitted propensities are clipped so that $1-\ehat(x)\ge c$.
\end{assumption}

\begin{remark}[Positivity is one-sided, and this is a real weakening]
\label{rem:positivity}
\textbf{Changed from all previous drafts}, which assumed the two-sided
$\ez(x)\in[c,1-c]$ while noting it was ``stronger than what is required only for
identification of the ATT''. For the ATT nothing requires a lower bound on
$\ez$:
\begin{enumerate}[nosep,label=(\roman*)]
  \item \emph{Identification.} $\E[Y(1)\mid A=1]$ is identified from treated
  units with no positivity at all; $\E[Y(0)\mid A=1]$ needs control units
  wherever treated units live, i.e.\ $\ez(x)<1$. If $\ez(x)=0$ there are no
  treated units at $x$ and $x$ does not enter the ATT.
  \item \emph{The weight.} $\wz=\ez/(1-\ez)$ is bounded iff $1-\ez\ge c$.
  \item \emph{The remainder.} By the display after \eqref{eq:bilinear}, only
  $1-\ez,1-\ehat\ge c$ is used.
  \item \emph{The margin} (Lemma~\ref{lem:margin}) needs positive control-arm
  mass in each cell, $p_x\{1-\ez(x)\}\ge p_xc>0$: again one-sided.
  \item \emph{The efficiency bound} $V=\pi^{-2}\E[\{A(Y-\mz-\theta_0)-(1-A)\wz(Y-\mz)\}^2]$
  is finite iff $\wz$ is bounded: one-sided.
\end{enumerate}
\textbf{Where the two-sided bound was actually used}, and how it is avoided: the
\emph{norm-to-risk} direction of the cross-entropy loss--norm link has constant
$U_c=\min\{1/(2c^2),1/(c(1-c))\}$, which diverges as $\ez\to0$ as well as
$\ez\to1$. That direction is only needed if the propensity is fitted by
cross-entropy. Fitting it by squared error on the binary $A$ --- which is what
within-leaf treated fractions compute anyway --- makes the link exact
(A\ref{ass:lossnorm}) and removes the dependence. If a cross-entropy solver is
used for structure search, the two-sided bound returns for that step only and
should be stated there rather than globally.
\end{remark}

\begin{assumption}[Second moments]\label{ass:moments}
$\E[Y^2\mid X=x]\le C_Y<\infty$ for every $x\in\cX$.
\end{assumption}

\begin{assumption}[Loss--norm link]\label{ass:lossnorm}
For squared error the excess risk \emph{equals} $\Ltwo{f-\gamma_0}^2$, so
$C_{\mathrm{ln}}=1$.
\end{assumption}

\begin{remark}[Constant correction]
Earlier drafts state $C_{\mathrm{ln}}=\tfrac12$ for squared error. The excess
risk equals the squared $L_2$ error exactly; $1$ is correct.
\end{remark}

\begin{assumption}[Global optimization]\label{ass:global}
$\that_j$ solves \eqref{eq:select} globally over $\cT_{\bar L}$, with
$\lambda_n\to0$.
\end{assumption}

\subsection{Standing conditions on the fitted trees}

These are conditions on the \emph{estimator}, not on the truth. They are stated
once and used throughout.

\begin{assumption}[Leaf-level clipping]\label{ass:clip}
Fitted propensities are clipped so that $1-\ehat(x)\ge c$.
\end{assumption}

\begin{remark}[What A\ref{ass:clip} is for, in plain terms]\label{rem:clip}
The fitted propensity in a leaf is the fraction of treated units in that leaf.
If a leaf happens to contain \emph{only} treated units then $\ehat=1$ exactly and
the weight $\ehat/(1-\ehat)$ divides by zero, so the estimator returns a
non-finite value. For a leaf holding $m$ observations the probability of this is
at most $(1-c)^m$ --- small, but not zero, and across thousands of simulation
replications it will occur. There are two remedies: require leaves large enough
that it is negligible ($m_n\gtrsim\log n$), or clip. Clipping is one line and
removes the failure mode outright, so we clip.
\end{remark}

\begin{assumption}[Minimum leaf mass]\label{ass:mass}
Each leaf of $\that_j$ contains at least $m_n$ observations, with
$\bar L\,m_n\lesssim n$.
\end{assumption}

\begin{remark}[What A\ref{ass:mass} is for, in plain terms]\label{rem:mass}
A leaf mean estimated from a leaf of probability mass $p_\ell$ has variance
$\sigma^2_\ell/(np_\ell)$, so summing over leaves,
\[
  \E\Ltwo{\mhat-\mz}^2 \;\asymp\; \frac1n\sum_{\ell}\frac{\sigma^2_\ell}{p_\ell}.
\]
If the leaves are balanced, $p_\ell\asymp1/\bar L$, this is $\asymp\bar L/n$ ---
the rate the entire argument uses. But a \emph{single} leaf of tiny mass makes
its $1/p_\ell$ term explode and the sum is then far larger than $\bar L/n$,
which breaks the product remainder. So we need $\min_\ell p_\ell\gtrsim1/\bar L$,
i.e.\ the leaves must not be too unbalanced. Operationally this is enforced as a
minimum leaf \emph{count} $m_n$; the side condition $\bar Lm_n\lesssim n$ merely
says one cannot demand more leaves times more points-per-leaf than there is data.
Neither this nor A\ref{ass:clip} is currently enforced or measured in the
software.
\end{remark}

\begin{assumption}[Leaf budget]\label{ass:budget}
$\card{\that_j}\le\bar L$, and in a growing regime $L_n=o(\sqrt n)$ with the
entropy written explicitly as $L_n\log(dJ)$.
\end{assumption}

\begin{remark}[Why $L_n=o(\sqrt n)$ is the right condition, and why the entropy
is $O(1)$ here]\label{rem:selfinfluence}
The self-influence of one observation on its own leaf mean is
$O(1/n_{\mathrm{leaf}})=O(L/n)$, which is $o(n^{-1/2})$ exactly when
$L=o(\sqrt n)$: every leaf holds far more than $\sqrt n$ observations, so no
single point can move its own prediction enough to matter. This is
\citet{chen2022}'s leave-one-out stability condition specialised to leaf means;
ours is the verification, not the idea.

On a finite product space with $d$ fixed there are only $d(J-1)$ candidate
splits, so $\log\card{\cT_{\bar L}}\asymp\bar L\log(dJ)=O(1)$ and the union bound
costs a \emph{constant}; $L_n\log n=o(\sqrt n)$ is then vacuous. \textbf{Part~II
must therefore not be motivated as escaping an $\exp(CL_n\log n)$ entropy cost}
--- our own finite-space assumption forbids that cost from existing.
\end{remark}

\subsection{Structural assumptions, by part}

\begin{definition}[Approximation error, weighted]\label{def:delta}
With $\Lw{f}^2=\E[\{1-\ez(X)\}f(X)^2]$,
\[
  \delta_e=\min_{\tau\in\cT_{\bar L}}\Ltwo{\Pi_\tau \ez-\ez},
  \qquad
  \delta_\mu=\min_{\tau\in\cT_{\bar L}}\Lw{\Pi_\tau\mz-\mz},
\]
where $\Pi_\tau$ is projection onto functions constant on $\tau$. The
$\mu$ error is \textbf{weighted}, per the weight in \eqref{eq:bilinear}. Under
A\ref{ass:causal} the weighted and unweighted norms are equivalent up to $1/c$,
so this affects constants and not rates --- but defining $\delta_\mu$ unweighted
and then substituting it into \eqref{eq:bilinear} is an error, not a cosmetic
choice.
\end{definition}

\begin{assumption}[Part~II: structural sparsity]\label{ass:sparsity}
$\bar L\ge\max(L^*_e,L^*_\mu)$; equivalently $\delta_e=\delta_\mu=0$.
\end{assumption}

\begin{remark}[This is sparsity, not correct specification]\label{rem:not-exact}
On a finite space \emph{every} function is piecewise constant --- the partition
into singletons always works, so $L^*_j\le M$ holds automatically. So
A\ref{ass:sparsity} asserts nothing about functional form; its entire content is
that $L^*_j$ is \emph{small enough to display}. It must not be described as
exactness or correct specification, and it is not vulnerable to ``you assumed the
truth is a step function''.
\end{remark}

\begin{remark}[Why trees rather than a sparse additive model]\label{rem:additive}
$L^*\le2^s$ always, and often far less: \emph{hierarchical} dependence, where the
active coordinate set differs by region, costs as few as $s+1$ leaves.
A\ref{ass:sparsity} is expensive precisely for \emph{additive} structure
($x_1+x_2+x_3$ needs a leaf per configuration, the worst case for trees) and
cheap precisely for region-varying active sets. That asymmetry, not flexibility,
is the reason to use trees here, and it is the answer to ``a sparse additive
model attains the same rate with none of this apparatus''. A concrete example
must accompany this remark.
\end{remark}

\begin{assumption}[Part~III: graded approximation]\label{ass:graded}
$\delta_e\cdot\delta_\mu=o(n^{-1/2})$. A sufficient form:
$\delta_j(L)\le C_jL^{-q_j}$ with $q_e+q_\mu>1$ and $\bar L\asymp\sqrt n$ up to
logarithmic factors.
\end{assumption}

%======================================================================
\section{Part I: cross-fitted inference with optimal-tree nuisances}
%======================================================================

\begin{assumption}[Rate condition, product form]\label{ass:rate}
\[
  \Ltwo{\ehat-\ez}\cdot\Lw{\mhat-\mz}\;=\;o_p(n^{-1/2}).
\]
In the continuous tier this is
$\bar\alpha_e/(2\bar\alpha_e+s_e)+\bar\alpha_\mu/(2\bar\alpha_\mu+s_\mu)>1/2$.
\end{assumption}

\begin{remark}[Why the product form, and why the per-nuisance form is unusable]
\label{rem:product}
\textbf{This is the highest-value correction in Part~I.} The per-nuisance
requirement $\Ltwo{\widehat\gamma_j-\gamma_{0,j}}=o_p(n^{-1/4})$ forces
$\bar\alpha_j>s_j/2$ for \emph{each} $j$. Since $s$ is a supremum of active-set
cardinalities it is an integer with no fractional relaxation, and
$\bar\alpha\le1$ for piecewise-constant trees, so $\bar\alpha>s/2$ forces
$\boldsymbol{s=1}$ --- locally univariate, excluding every smooth two-way
interaction. The product form requires only that the two ratios sum past
$1/2$, so one nuisance may be poorly estimated when the other is well estimated;
it admits $s\le3$ for one nuisance, and it decouples Part~I from Part~II. The
product form is already stated in the current draft and then abandoned three
times in favour of the per-nuisance version.
\end{remark}

\begin{theorem}[Cross-fitted ATT with optimal-tree nuisances]\label{thm:crossfit}
Under A\ref{ass:causal}--A\ref{ass:lossnorm}, A\ref{ass:rate}, and $K$-fold
cross-fitting with $K$ fixed, the estimator \eqref{eq:att-estimator} computed
with fold-specific nuisances satisfies
\[
  \sqrt n\,(\thetahat_{\mathrm{cf}}-\theta_0)\;\rightsquigarrow\;N(0,V),
  \qquad
  V=\pi^{-2}\,\E\Bigl[\bigl\{A(Y-\mz-\theta_0)-(1-A)\wz(Y-\mz)\bigr\}^2\Bigr].
\]
\end{theorem}

\begin{remark}[Rung two of the ladder: $K$ interpretable trees]
Theorem~\ref{thm:crossfit} imposes no leaf budget and no sparsity assumption, and
it yields $K$ trees per nuisance rather than one. Those trees are individually
readable and auditable fold by fold, which is already a gain over a black-box
nuisance; what it does not give is a single displayed model. It is also the
honest answer to ``how would you relax A\ref{ass:sparsity}''.
\end{remark}

%======================================================================
\section{Part II: the two-tree estimator, full sample, no splitting}
%======================================================================

Structure \emph{and} leaf values come from all $n$. No folds, no sample
splitting, no cross-fitting, and no intersection or averaging over multiple fits.

\begin{lemma}[Sufficiency and invariance]\label{lem:invariance}
Under A\ref{ass:sparsity}, for every $\tau\in\cS_j$ the induced nuisance equals
$\gamma_{0,j}$ on the support, and hence the induced estimand equals $\theta_0$
for every pair $(\tau_e,\tau_\mu)\in\cS_e\times\cS_\mu$.
\end{lemma}

\begin{remark}[This is where the novelty sits]
Structure is chosen by data-adaptive optimisation on the same sample used for
inference, so naive post-selection inference would not generally be valid. Here
it is, and the reason is specific and checkable: by Lemma~\ref{lem:invariance}
selection cannot move the target. Note the payoff is for the
\emph{distributional} claim only --- it does not license a bias or MSE claim.
\end{remark}

\begin{lemma}[The margin is free under global optimisation]\label{lem:margin}
Under A\ref{ass:causal}--A\ref{ass:sparsity},
$\Delta_j:=\min_{\tau\in\cT_{\bar L}\setminus\cS_j}\{R^{(j)}(\tau)-R^{(j)}(\tau^*_j)\}>0$.
\end{lemma}

\begin{remark}[Global optimisation supplies the margin; greedy does not]
\label{rem:greedy}
Lemma~\ref{lem:margin} is automatic, not an assumption, because $\cT_{\bar L}$ is
finite and any $\tau$ failing to refine $\tau^*_j$ has a leaf containing two
distinct values of $\gamma_{0,j}$, both of positive probability. Under greedy
recursive splitting no such margin exists: if $\ez$ depends on $X$ only through
$\mathbf1\{x_1=x_2\}$, the marginal risk reduction of \emph{either} correct first
split is exactly zero, so $\Delta=0$. The result therefore genuinely requires an
exact solver. This is one place where using optimal rather than CART-style trees
does mathematical rather than rhetorical work, and it strengthens Part~I's
novelty claim as well.
\end{remark}

\begin{lemma}[Selection lands in the sufficient class]\label{lem:selection}
Under A\ref{ass:causal}--A\ref{ass:global} and A\ref{ass:sparsity},
$\PP(\that_j\in\cS_j)\to1$; if $Y$ is bounded,
$\PP(\that_j\notin\cS_j)\le\card{\cT_{\bar L}}e^{-c_1n\Delta_j^2}$.
\end{lemma}

\begin{lemma}[Conditional finite-dimensionality]\label{lem:findim}
For fixed $\tau$, the refitted nuisance is determined by at most $\bar L$ real
numbers per nuisance, each a within-leaf mean converging at $O_p(n^{-1/2})$.
\end{lemma}

\begin{lemma}[Cross term via leaf-wise conditional variance]\label{lem:cross}
With $\that_e\ne\that_\mu$ carried separately, the term linear in
$\ehat-\ez$ satisfies $\abs{T}=O_p(\sqrt{\bar L}\log n/n)=o_p(n^{-1/2})$.
\end{lemma}

\begin{remark}[Cauchy--Schwarz is fatal here, not merely lossy]\label{rem:cs}
Bounding $T$ by Cauchy--Schwarz gives
$\Ltwo{\ehat-\ez}\cdot O_p(\sqrt{\bar L/n})$ whose second factor is the
\emph{irreducible outcome standard deviation} $O_p(1)$, not $\delta_\mu$ ---
yielding $O_p(\sqrt{\bar L\log n/n})$, which is not $o(n^{-1/2})$ and silently
reintroduces $s=1$. The leaf-wise variance argument instead centres the weight
within each outcome leaf, so
$\sum_{i\in\ell,A_i=0}(\what-\bar w_\ell)=0$ removes the leading piece and leaves
a conditionally mean-zero sum. It requires conditionally mean-zero residuals
(\emph{not} independence across leaves) and measurable coefficients, and it has
large slack. Note the cross terms do \emph{not} vanish under
A\ref{ass:sparsity}: exactness only shrinks the input.
\end{remark}

\begin{theorem}[Full-sample two-tree CLT]\label{thm:main}
Under A\ref{ass:causal}--A\ref{ass:budget} and A\ref{ass:sparsity}, with
$\that_e,\that_\mu$ selected by \eqref{eq:select} on all $n$ and leaf values
refit on all $n$,
\[
  \sqrt n\,(\thetahat-\theta_0)\;\rightsquigarrow\;N(0,V),
\]
with $V$ as in Theorem~\ref{thm:crossfit}.
\end{theorem}

\begin{corollary}[Variance estimation and confidence intervals]\label{cor:variance}
Let
\begin{equation}\label{eq:vhat}
  \Vhat \;=\; \frac{1}{n\,\pihat^{2}}\sum_{i=1}^n
  \Bigl[\,A_i\bigl\{Y_i-\mhat(X_i)-\thetahat\bigr\}
   \;-\;(1-A_i)\,\what(X_i)\bigl\{Y_i-\mhat(X_i)\bigr\}\Bigr]^2 .
\end{equation}
Then $\Vhat\to_p V$ and
$\thetahat\pm z_{1-\alpha/2}\sqrt{\Vhat/n}$ has asymptotic coverage $1-\alpha$.
Estimating $\pi$ is free: $(\pihat-\pi)\Pn\psi_0=O_p(n^{-1})$.
\end{corollary}

\begin{remark}[This formula did not previously exist]
Earlier drafts assert that a consistent $\Vhat$ exists without giving one, and no
leave-one-out or jackknife variance estimator appears anywhere in the software;
what ships is a cross-fitted twin's standard error inside a bias-aware interval,
which validates a different construction. Equation~\eqref{eq:vhat} is the
ordinary empirical influence-function variance, and it is consistent here
\emph{because} of Lemma~\ref{lem:findim}: conditional on a sufficient partition
the problem is finite-dimensional, so no sample-splitting correction is needed.
\end{remark}

\begin{remark}[The inference problem here is finite-dimensional; say so]
\label{rem:findim-objection}
With $M$ fixed, $P$ is determined by finitely many parameters and the ATT is a
smooth function of them, so Theorem~\ref{thm:main} is \emph{not} a nonparametric
efficiency result and must not be sold as one. Conditional on a sufficient
partition the estimator is an ordinary smooth $Z$-estimator in at most
$2\bar L+1$ parameters --- which is precisely why no sample splitting is
required. Stating this disarms the objection; omitting it invites it.
\end{remark}

\begin{remark}[Why two partitions, and why a shared one degenerates]
\label{rem:shared}
Computing both leaf means on the intersection partition
$\that_e\wedge\that_\mu$ would make $\what$ constant within each outcome leaf and
drive $T$ to exactly zero. It is tempting and it must be resisted: the same
mechanism collapses the estimator to a stratified difference in means in which
\textbf{the propensity plays no role whatsoever} --- $\what$ would influence
$\thetahat$ only through the partition on which $\mhat$ is computed, never
through its values. One of the two displayed trees would be decorative. Since
Lemma~\ref{lem:cross} controls $T$ with large slack, there is no need to buy it
off. The nonzero cross term is what keeps both trees load-bearing.
\end{remark}

%======================================================================
\section{Part III: honest inference when sparsity fails}
%======================================================================

Under A\ref{ass:graded} in place of A\ref{ass:sparsity},
Lemma~\ref{lem:invariance} fails and the bilinear remainder
\eqref{eq:bilinear} no longer vanishes.

\begin{lemma}[Do not subtract; widen]\label{lem:no-subtract}
$\thetahat-\delta_{\mathrm{fid}}=\thetahat_{\mathrm{cf}}$ identically. Debiasing
by the fidelity diagnostic therefore returns the cross-fitted estimator exactly,
discarding the interpretable point estimate and gaining nothing.
\end{lemma}

\begin{theorem}[Honest interval outside the class]\label{thm:honest}
Under A\ref{ass:causal}--A\ref{ass:budget} and A\ref{ass:graded},
\[
  \thetahat \;\pm\; \Bigl(z_{1-\alpha/2}\sqrt{\Vhat/n}
    \;+\; \frac{\delta_e\delta_\mu}{\pi c}\Bigr)
\]
has asymptotic coverage at least $1-\alpha$ for $\theta_0$.
\end{theorem}

\begin{corollary}[The price, and where it plateaus]\label{cor:width}
The width is $\asymp\max\{n^{-1/2},\,\delta_e\delta_\mu\}$: the interval behaves
like the efficient one while the approximation error is negligible, and
plateaus at $\delta_e\delta_\mu$ once it is not.
\end{corollary}

\begin{proposition}[$\delta$ is estimable when $M\ll n$]\label{prop:estimable}
On a finite space the saturated cell-wise fit is available, so
$\delta_j^2$ is identified as the excess risk of the best $\bar L$-leaf tree
relative to it, and a held-out estimator $\dhat_j$ satisfies
$\dhat_j^2-\delta_j^2=O_p(\sqrt{M/n})$. Substituting $\dhat_e\dhat_\mu$ into
Theorem~\ref{thm:honest}, with an allowance for its own sampling error, gives a
\emph{checkable} bias-aware interval requiring no assumed decay rate.
\end{proposition}

\begin{remark}[Scope of Proposition~\ref{prop:estimable}, honestly stated]
\label{rem:scope}
The saturated fit must be estimable, so $M\ll n$. Because $\mz$ is fit on
controls only the effective sample is about $n/2$, and the binding constraint is
the \emph{smallest} cell rather than the average. At roughly $30$ observations
per cell:

\begin{center}
\begin{tabular}{@{}rrr@{}}
\toprule
$n$ & $M$ & binary covariates \\
\midrule
$1{,}000$ & $\lesssim 17\text{--}33$ & $d\le4\text{--}5$ \\
$10{,}000$ & $\lesssim 170\text{--}330$ & $d\le7\text{--}8$ \\
$100{,}000$ & $\lesssim 1{,}700\text{--}3{,}300$ & $d\le10\text{--}11$ \\
\bottomrule
\end{tabular}
\end{center}

So $d$ grows like $\log_2 n-5$: this is a genuinely restrictive regime and must
be reported as such. Note it constrains \textbf{only}
Proposition~\ref{prop:estimable}. Theorem~\ref{thm:main} and
Theorem~\ref{thm:honest} leave $d$ unconstrained --- there the binding constraint
is interpretability itself, since a tree with many leaves defeats the purpose.
\end{remark}

\begin{remark}[No decay rate is implied by the domain]\label{rem:parity}
A\ref{ass:graded} cannot be derived from finiteness. Parity on $d$ binary
coordinates has $\delta(L)$ bounded away from zero for every axis-aligned tree
with $L<2^{d-1}$ leaves and then drops to zero, so decay can be a step function
of $L$. A\ref{ass:graded} is the discrete analogue of assuming a H\"older or
Besov ball, which is equally ``assuming the conclusion'' via Jackson-type
theorems and which nobody objects to. A sufficient condition should be supplied
as a lemma: weak-$\ell_r$ decay of tree/Haar coefficients gives $q=1/r-1/2$.
Low-order-interaction or ANOVA decay must \textbf{not} be used --- such functions
are well approximated by \emph{additive} models, which is trees' worst case
(Remark~\ref{rem:additive}).
\end{remark}

%======================================================================
\section{Changes from previous drafts}\label{sec:changes}
%======================================================================

\begin{enumerate}[label=(\arabic*)]
  \item \textbf{Positivity weakened to one-sided} (A\ref{ass:causal},
  Remark~\ref{rem:positivity}). Was $\ez\in[c,1-c]$; now $1-\ez\ge c$. The
  two-sided version was used only by the norm-to-risk direction of the
  cross-entropy link, which squared-error fitting of the propensity avoids.
  \item \textbf{Rate condition in product form} (A\ref{ass:rate},
  Remark~\ref{rem:product}). The per-nuisance form forces $s=1$.
  \item \textbf{$C_{\mathrm{ln}}=1$}, not $\tfrac12$, for squared error.
  \item \textbf{Sign of the control-arm term} corrected in
  \eqref{eq:att-estimator} (Remark~\ref{rem:sign}).
  \item \textbf{$\delta_\mu$ defined on the $(1-\ez)$-weighted projection}
  (Definition~\ref{def:delta}).
  \item \textbf{An explicit variance estimator} \eqref{eq:vhat}, which did not
  previously exist.
  \item \textbf{Five symbol collisions resolved} (Section~\ref{sec:collisions}),
  including the $\delta$ / $\delta_{\mathrm{fid}}$ conflation carried by the
  software (Remark~\ref{rem:two-deltas}).
  \item \textbf{Rashomon sets and cross-fold intersection removed} from the main
  line. The displayed-tree claim is about \emph{count}, which the full-sample
  estimator delivers directly; the intersection additionally forced a coarse
  discretisation that starved structure selection.
\end{enumerate}

\paragraph{Open, and needed before drafting proofs.}
Proofs of Theorem~\ref{thm:main} (assembly), Lemma~\ref{lem:cross} (the
leaf-wise variance argument in full), Theorem~\ref{thm:honest}, and
Proposition~\ref{prop:estimable} (the $O_p(\sqrt{M/n})$ rate, and how
$\dhat$'s sampling error enters the interval). In Part~I,
the near-ERM oracle inequality currently bounds terms by
$O_p(\sqrt{L_n\log n/n})$ and then squares them inside an additive risk
decomposition, which is not a valid step and understates the rate by a square;
it needs the localised sub-root argument.

\begin{thebibliography}{9}
\bibitem[Chen et al.(2022)]{chen2022} Chen, Q., Syrgkanis, V., Austern, M.
  (2022). Debiased machine learning without sample splitting.
  arXiv:2206.01825.
\bibitem[Hahn(1998)]{hahn1998} Hahn, J. (1998). On the role of the propensity
  score in efficient semiparametric estimation of average treatment effects.
  \emph{Econometrica} 66, 315--331.
\bibitem[Kennedy(2024)]{kennedy2024} Kennedy, E. H. (2024). Semiparametric
  doubly robust targeted double machine learning: a review.
\bibitem[Robins \& Rotnitzky(1995)]{robins1995} Robins, J. M., Rotnitzky, A.
  (1995). Semiparametric efficiency in multivariate regression models with
  missing data. \emph{JASA} 90, 122--129.
\bibitem[van der Laan \& Robins(2003)]{vanderlaan2003} van der Laan, M. J.,
  Robins, J. M. (2003). \emph{Unified Methods for Censored Longitudinal Data
  and Causality}. Springer.
\end{thebibliography}

\end{document}
